%  LaTeX - Sistema de Asignación de Salones ISC
\documentclass[algorithms,article,submit,pdftex]{Definitions/mdpi} 

\firstpage{1} 
\makeatletter 
\setcounter{page}{\@firstpage} 
\makeatother
\pubvolume{1}
\issuenum{1}
\articlenumber{0}
\pubyear{2026}
\copyrightyear{2025}

\Title{Sistema de Asignación de Salones ISC: Optimización Inteligente de Espacios Académicos}

\Author{Anónimo}
\AuthorNames{Anónimo}

\address{Instituto Tecnológico de Ciudad Madero, Tecnológico Nacional de México}

\corres{Correspondencia: contacto@ejemplo.com}

\abstract{Sistema inteligente de optimización para la asignación de 680 clases a 21 salones en el programa de Ingeniería en Sistemas Computacionales del Instituto Tecnológico de Ciudad Madero. El sistema implementa cuatro algoritmos de optimización: heurística del profesor, Greedy con Hill Climbing, Machine Learning (Random Forest) y Algoritmo Genético. Se presenta un mecanismo de pre-asignación que garantiza el 100\% de cumplimiento de las preferencias prioritarias (Prioridad 1) mientras optimiza tres objetivos: minimizar movimientos de profesores entre salones, reducir cambios de piso y disminuir la distancia total recorrida. Los resultados experimentales de 90 ejecuciones demuestran que el enfoque Greedy+Hill Climbing logra el mejor rendimiento general con 314 movimientos (mejora del 12\% sobre la línea base), 206 cambios de piso (reducción del 28\%) y resultados consistentes. La validación estadística mediante ANOVA, pruebas post-hoc de Tukey HSD y análisis de tamaño de efecto de Cohen confirma diferencias significativas entre todos los algoritmos. El sistema está implementado en Python con documentación completa y está disponible como software de código abierto.}

\keyword{asignación de salones; optimización combinatoria; algoritmos híbridos; hill climbing; machine learning; algoritmo genético; programación académica; satisfacción de restricciones}

\begin{document}

\section{Sistema de Asignación de Salones ISC}

\textbf{Optimización Inteligente de Espacios Académicos}

\textbf{Jesús Olvera}

Ingeniería en Sistemas Computacionales

Instituto Tecnológico de Ciudad Madero


\subsection{Descripción General}

Sistema inteligente de optimización para la asignación de salones en el programa de Ingeniería en Sistemas Computacionales.

\textbf{Objetivos:}

\begin{itemize}
\item  Minimizar movimientos de profesores
\begin{itemize}
\item  Reducir cambios de piso
\begin{itemize}
\item  Optimizar distancias recorridas
\begin{itemize}
\item  Garantizar cumplimiento de preferencias prioritarias

\subsection{ Sistema de Prioridades Jerárquico}

\textbf{Tres niveles de prioridad:}

\begin{enumerate}
\item \textbf{PRIORIDAD 1 (Hard Constraint):} Preferencias de Profesores
\begin{itemize}
\item Cumplimiento: \textbf{100% garantizado}
\begin{itemize}
\item Implementación: Pre-asignación forzada
\begin{itemize}
\item Protección: Clases inmutables durante optimización
\begin{enumerate}
\item \textbf{PRIORIDAD 2 (Soft Constraint):} Consistencia de Grupos
\begin{itemize}
\item Mantener grupos en el mismo salón
\begin{enumerate}
\item \textbf{PRIORIDAD 3 (Soft Constraint):} Primer Semestre
\begin{itemize}
\item Asignar grupos 15xx a salones específicos

\subsection{ Algoritmos de Optimización}

\textbf{4 optimizadores diferentes:}


\subsection{Arquitectura del Sistema}

\begin{verbatim}
Sistema-Salones-ISC/
 configurador_materias.py      # Interfaz gráfica
 pre_asignar_p1.py             # Pre-asignación P1
 optimizador_greedy.py         # Greedy + Hill Climbing
 optimizador_ml.py             # Machine Learning
 optimizador_genetico.py       # Algoritmo Genético
 corregir_prioridades.py       # Corrección post-opt
 ejecutar_todos.py             # Script maestro
 generar_comparativa_completa.py  # Reportes
 utils_restricciones.py        # Validación
 datos_estructurados/          # Datos I/O
 comparativas/                 # Resultados
\end{verbatim}


\subsection{Flujo de Ejecución}

\begin{verbatim}
graph TD
    A[Horario Inicial] --> B[Pre-asignación P1]
    B --> C[00_Horario_PreAsignado_P1.csv]
    C --> D[Optimizador Greedy]
    C --> E[Optimizador ML]
    C --> F[Optimizador Genético]
    D --> G[Corrección Post-Opt]
    E --> H[Corrección Post-Opt]
    F --> I[Corrección Post-Opt]
    G --> J[Comparativas y Gráficos]
    H --> J
    I --> J
    J --> K[Reportes Finales]
\end{verbatim}


\subsection{Métricas de Optimización}

\textbf{Función Objetivo} minimiza:

\begin{itemize}
\item \textbf{Movimientos de profesores:} Cambios de salón
\begin{itemize}
\item \textbf{Cambios de piso:} Subir/bajar escaleras
\begin{itemize}
\item \textbf{Distancia total:} Recorrido acumulado
\begin{itemize}
\item \textbf{Penalizaciones:} Violaciones de restricciones soft

\subsection{Resultados Típicos}


\section{Contexto del Problema}

\textbf{Asignación Óptima de Salones}


\subsection{1. Introducción - Planteamiento del Problema}

La asignación de salones en instituciones educativas es un \textbf{problema de optimización combinatoria NP-completo} que surge de la necesidad de distribuir eficientemente los espacios físicos disponibles entre las diferentes actividades académicas programadas.

\textbf{Complejidad:}

\begin{itemize}
\item Múltiples restricciones simultáneas
\begin{itemize}
\item Objetivos conflictivos
\begin{itemize}
\item Espacio de búsqueda exponencial

\subsection{Contexto Institucional - ITCM}

\textbf{Recursos Disponibles:}

\begin{itemize}
\item \textbf{Salones de Teoría:} 13 aulas (FF1-FF7 en planta baja, FF8-FFD en planta alta)
\begin{itemize}
\item \textbf{Laboratorios:} 8 espacios especializados (LBD, LBD2, LCA, LCG1, LCG2, LIA, LR, LSO)
\begin{itemize}
\item \textbf{Salones Inválidos:} 5 espacios no utilizables (AV1, AV2, AV4, AV5, E11)
\textbf{Demanda:}

\begin{itemize}
\item \textbf{Clases totales:} ~680 sesiones semanales
\begin{itemize}
\item \textbf{Materias:} 37 diferentes
\begin{itemize}
\item \textbf{Grupos:} Distribuidos en 9 semestres
\begin{itemize}
\item \textbf{Profesores:} ~30 docentes con preferencias específicas

\subsection{Complejidad del Problema}

El espacio de búsqueda para este problema es:

\begin{equation}
|\Omega| = m^n = 21^{680} \approx 10^{900}
\end{equation}

Donde:

\begin{itemize}
\item $m = 21$ salones disponibles
\begin{itemize}
\item $n = 680$ clases a asignar
\textbf{Perspectiva:}

\begin{itemize}
\item \textbf{Átomos en el universo observable:} $\approx 10^{80}$
\begin{itemize}
\item \textbf{Combinaciones posibles:} $\approx 10^{900}$
\begin{itemize}
\item \textbf{Relación:} $10^{820}$ veces más combinaciones que átomos
 Esto hace \textbf{imposible} la búsqueda exhaustiva


\subsection{2. Formulación Matemática - Conjuntos Básicos}

\begin{equation}
\begin{align}
C &= \{c_1, c_2, ..., c_n\} &&\text{Conjunto de clases} \\
S &= \{s_1, s_2, ..., s_m\} &&\text{Conjunto de salones} \\
P &= \{p_1, p_2, ..., p_k\} &&\text{Conjunto de profesores} \\
D &= \{\text{Lunes, Martes, ..., Viernes}\} &&\text{Dias de la semana} \\
H &= \{0700, 0800, ..., 2100\} &&\text{Bloques horarios}
\end{align}
\end{equation}


\textbf{Donde:}

\begin{itemize}
\item $n = 680$ (número total de clases)
\begin{itemize}
\item $m = 21$ (número de salones disponibles)
\begin{itemize}
\item $k \approx 30$ (número de profesores)
\begin{itemize}
\item $\{c_1, c_2, ..., c_n\}$ denota el conjunto de todas las clases
\begin{itemize}
\item $\in$ significa "pertenece a" o "es elemento de"

\subsection{Atributos de Clases}

Para cada clase $c_i \in C$:

\begin{equation}
\begin{align}
dia(c_i) &\in D &&\text{Dia de la semana} \\
hora(c_i) &\in H &&\text{Bloque horario} \\
materia(c_i) &\in M &&\text{Materia} \\
grupo(c_i) &\in G &&\text{Grupo} \\
profesor(c_i) &\in P &&\text{Profesor asignado} \\
tipo(c_i) &\in \{\text{Teoria, Laboratorio}\} &&\text{Tipo de clase} \\
estudiantes(c_i) &\in \mathbb{N} &&\text{Numero de estudiantes}
\end{align}
\end{equation}


\textbf{Donde:}

\begin{itemize}
\item $c_i$ = clase individual (ejemplo: $c_1, c_2, ..., c_{680}$)
\begin{itemize}
\item $dia(c_i)$ = función que retorna el día de la clase $c_i$
\begin{itemize}
\item $\mathbb{N}$ = números naturales (1, 2, 3, ...)
\begin{itemize}
\item $M$ = conjunto de todas las materias
\begin{itemize}
\item $G$ = conjunto de todos los grupos

\subsection{Atributos de Salones}

Para cada salón $s_j \in S$:

\begin{equation}
\begin{align}
capacidad(s_j) &\in \mathbb{N} &&\text{Capacidad maxima} \\
tipo(s_j) &\in \{\text{Teoria, Laboratorio}\} &&\text{Tipo de salon} \\
piso(s_j) &\in \{0, 1\} &&\text{Planta baja o alta} \\
ubicacion(s_j) &\in \mathbb{R}^2 &&\text{Coordenadas fisicas}
\end{align}
\end{equation}


\textbf{Donde:}

\begin{itemize}
\item $s_j$ = salón individual (ejemplo: FF1, FF2, LAB1, ...)
\begin{itemize}
\item $\mathbb{R}^2$ = plano de coordenadas reales (x, y)
\begin{itemize}
\item $\{0, 1\}$ = conjunto con dos elementos (0 = planta baja, 1 = planta alta)

\textbf{Variable de Decisión:}

\begin{equation}
A: C \rightarrow S
\end{equation}

\textbf{Donde:}

\begin{itemize}
\item $A$ = función de asignación
\begin{itemize}
\item $C \rightarrow S$ = función que mapea de clases a salones
\begin{itemize}
\item $A(c_i) = s_j$ significa que la clase $c_i$ se imparte en el salón $s_j$
\begin{itemize}
\item $\rightarrow$ = símbolo de función ("mapea a")

\subsection{Restricciones Duras (Hard Constraints)}

Estas restricciones \textbf{DEBEN} cumplirse para que la solución sea factible.

\textbf{R1. No Conflictos Temporales:}

\begin{equation}
\forall c_i, c_j \in C: \left(dia(c_i) = dia(c_j) \land hora(c_i) = hora(c_j) \land i \neq j\right) \Rightarrow A(c_i) \neq A(c_j)
\end{equation}

\textbf{Donde:}

\begin{itemize}
\item $\forall$ = "para todo" (cuantificador universal)
\begin{itemize}
\item $\land$ = "y" lógico (AND)
\begin{itemize}
\item $\Rightarrow$ = "implica que" (implicación lógica)
\begin{itemize}
\item $\neq$ = "diferente de" (no igual)

\textbf{R2. Capacidad Suficiente:}

\begin{equation}
\forall c_i \in C: estudiantes(c_i) \leq capacidad(A(c_i))
\end{equation}

\textbf{Donde:}

\begin{itemize}
\item $\leq$ = "menor o igual que"
\begin{itemize}
\item $estudiantes(c_i)$ = número de estudiantes en la clase $c_i$
\begin{itemize}
\item $capacidad(A(c_i))$ = capacidad del salón asignado a $c_i$

\subsection{Restricciones Duras (continuación)}

\textbf{R3. Tipo Correcto:}

\begin{equation}
\forall c_i \in C: tipo(c_i) = tipo(A(c_i))
\end{equation}

\textbf{Donde:}

\begin{itemize}
\item $tipo(c_i)$ = tipo de la clase (Teoría o Laboratorio)
\begin{itemize}
\item $tipo(A(c_i))$ = tipo del salón asignado
\begin{itemize}
\item Debe haber coincidencia exacta
\textit{Ejemplo: Una clase de laboratorio debe estar en un laboratorio}


\textbf{R4. Salones Válidos:}

\begin{equation}
\forall c_i \in C: A(c_i) \notin S_{invalidos}
\end{equation}

\textbf{Donde:}

\begin{itemize}
\item $\notin$ = "no pertenece a"
\begin{itemize}
\item $S_{invalidos} = \{AV1, AV2, AV4, AV5, E11\}$ = salones no utilizables

\textbf{R5. Preferencias Prioritarias (PRIORIDAD 1):}

\begin{equation}
\forall c_i \in P_1: A(c_i) = pref(c_i)
\end{equation}

\textbf{Donde:}

\begin{itemize}
\item $P_1$ = conjunto de clases con PRIORIDAD 1
\begin{itemize}
\item $pref(c_i)$ = salón preferido para la clase $c_i$
\begin{itemize}
\item Esta restricción es OBLIGATORIA (100% cumplimiento)

\subsection{Teorema 1: Factibilidad}

\textbf{Teorema 1 (Factibilidad):}

\textit{Una solución es factible si y solo si satisface todas las restricciones duras R1-R5.}

\textbf{Demostración:}

\begin{itemize}
\item $(\Rightarrow)$ Por definición de factibilidad
\begin{itemize}
\item $(\Leftarrow)$ Si se satisfacen R1-R5, no hay violaciones de restricciones obligatorias, por lo tanto la solución es válida

\subsection{Restricciones Suaves (Soft Constraints)}

Estas restricciones son \textbf{deseables} pero no obligatorias.

\textbf{S1. Consistencia de Grupo (PRIORIDAD 2):}

\begin{equation}
minimize \quad \left| \{A(c) : c \in C_g\} \right|
\end{equation}

\textbf{Donde:}

\begin{itemize}
\item $|\cdot|$ = cardinalidad (tamaño del conjunto)
\begin{itemize}
\item $\{A(c) : c \in C_g\}$ = conjunto de salones usados por el grupo $g$
\begin{itemize}
\item $C_g$ = clases del grupo $g$
\begin{itemize}
\item Objetivo: minimizar salones diferentes por grupo
\textit{Minimizar el número de salones diferentes usados por cada grupo}


\textbf{S2. Primer Semestre (PRIORIDAD 3):}

\begin{equation}
maximize \sum_{g \in G_{15}} \sum_{c \in C_g} \mathbb{1}[A(c) = salon_asignado(g)]
\end{equation}

\textbf{Donde:}

\begin{itemize}
\item $\sum$ = sumatoria (suma de todos los elementos)
\begin{itemize}
\item $G_{15}$ = grupos de primer semestre (15xx)
\begin{itemize}
\item $\mathbb{1}[\cdot]$ = función indicadora (1 si verdadero, 0 si falso)
\begin{itemize}
\item $salon\_asignado(g)$ = salón designado para el grupo $g$

\subsection{Restricciones Suaves (continuación)}

\textbf{S3. Minimizar Movimientos:}

Para cada profesor $p \in P$:

\begin{equation}
minimize \sum_{p \in P} \left| \{A(c) : c \in C_p\} \right|
\end{equation}

\textbf{Donde:}

\begin{itemize}
\item $C_p$ = conjunto de clases del profesor $p$
\begin{itemize}
\item $\{A(c) : c \in C_p\}$ = salones diferentes usados por el profesor $p$
\begin{itemize}
\item Objetivo: reducir salones diferentes por profesor
\textbf{Objetivo:} Reducir la cantidad de salones diferentes que usa cada profesor durante el día/semana


\subsection{Función Objetivo Completa}

\begin{equation}
\begin{align}
f(A) = &\ w_1 \cdot  movimientos(A) + w_2 \cdot  cambios_piso(A) \\
       &+ w_3 \cdot  distancia(A) + w_4 \cdot  penalizacion_invalidos(A) \\
       &+ w_5 \cdot  penalizacion_conflictos(A) + w_6 \cdot  penalizacion_tipo(A) \\
       &+ w_7 \cdot  penalizacion_P2(A) + w_8 \cdot  penalizacion_P3(A)
\end{align}
\end{equation}

\textbf{Donde:}

\begin{itemize}
\item $f(A)$ = energía o costo total de la asignación $A$
\begin{itemize}
\item $w_i$ = peso del componente $i$ (importancia relativa)
\begin{itemize}
\item $\cdot$ = multiplicación
\begin{itemize}
\item Cada componente mide un aspecto diferente de la calidad
\textbf{Objetivo:} $minimize\ f(A)$ sujeto a restricciones R1-R5


\subsection{Componente: Movimientos de Profesores}

\begin{equation}
movimientos(A) = \sum_{p \in P} \max\left(0, \left| \{A(c) : c \in C_p\} \right| - 1\right)
\end{equation}

\textbf{Donde:}

\begin{itemize}
\item $\max(a, b)$ = máximo entre $a$ y $b$
\begin{itemize}
\item $|\{A(c) : c \in C_p\}|$ = número de salones diferentes del profesor $p$
\begin{itemize}
\item Si usa $k$ salones, hace $k-1$ movimientos
\begin{itemize}
\item $\max(0, \cdot)$ asegura que el valor no sea negativo

\textbf{Ejemplo Numérico:}

\begin{itemize}
\item Profesor tiene clases en: FF1, FF2, FF1, FF3, FF2
\begin{itemize}
\item Salones únicos: {FF1, FF2, FF3}
\begin{itemize}
\item Movimientos = |{FF1, FF2, FF3}| - 1 = \textbf{2 movimientos}

\subsection{Componente: Cambios de Piso}

Sea $C_p^{ordenado} = [c_1, c_2, ..., c_k]$ las clases del profesor $p$ ordenadas por $(dia, hora)$:

\begin{equation}
cambios_piso(A) = \sum_{p \in P} \sum_{i=1}^{|C_p|-1} \mathbb{1}[piso(A(c_i)) \neq piso(A(c_{i+1}))]
\end{equation}


\textbf{Donde:}

\begin{itemize}
\item $C_p^{ordenado}$ = clases del profesor ordenadas cronológicamente
\begin{itemize}
\item $\sum_{i=1}^{|C_p|-1}$ = suma desde la primera hasta la penúltima clase
\begin{itemize}
\item $c_i, c_{i+1}$ = clases consecutivas en el tiempo
\begin{itemize}
\item $piso(A(c_i))$ = piso del salón asignado a la clase $c_i$
Donde $\mathbb{1}[\cdot]$ es la función indicadora:

\begin{equation}
\mathbb{1}[condicion] = \begin{cases}
1 & \text{si condicion es verdadera} \\
0 & \text{si condicion es falsa}
\end{cases}
\end{equation}


\subsection{Componente: Distancia Total}

Función de distancia entre salones:

\begin{equation}
d(s_i, s_j) = \begin{cases}
0 & \text{si } s_i = s_j \\
1 & \text{si } piso(s_i) = piso(s_j) \land s_i \neq s_j \\
10 & \text{si } piso(s_i) \neq piso(s_j)
\end{cases}
\end{equation}


\textbf{Donde:}

\begin{itemize}
\item $d(s_i, s_j)$ = distancia entre salón $s_i$ y salón $s_j$
\begin{itemize}
\item Mismo salón: distancia 0
\begin{itemize}
\item Mismo piso, diferente salón: distancia 1
\begin{itemize}
\item Diferente piso: distancia 10 (penalización por subir/bajar escaleras)

\textbf{Distancia total:}

\begin{equation}
distancia(A) = \sum_{p \in P} \sum_{i=1}^{|C_p|-1} d(A(c_i), A(c_{i+1}))
\end{equation}

\textbf{Donde:}

\begin{itemize}
\item Se suma la distancia entre salones de clases consecutivas
\begin{itemize}
\item Para cada profesor, se acumula la distancia total recorrida

\subsection{Componente: Penalizaciones}

\begin{equation}
\begin{align}
penalizacion_invalidos(A) &= 1000 \times \sum_{c \in C} \mathbb{1}[A(c) \in S_{invalidos}] \\
penalizacion_conflictos(A) &= 500 \times |\{(c_i, c_j) : conflicto_temporal(c_i, c_j, A)\}| \\
penalizacion_tipo(A) &= 300 \times \sum_{c \in C} \mathbb{1}[tipo(c) \neq tipo(A(c))] \\
penalizacion_P2(A) &= 50 \times \sum_{c \in P_2} \mathbb{1}[A(c) \neq pref(c)] \\
penalizacion_P3(A) &= 25 \times \sum_{c \in P_3} \mathbb{1}[A(c) \neq pref(c)]
\end{align}
\end{equation}


\textbf{Donde:}

\begin{itemize}
\item $\times$ = multiplicación
\begin{itemize}
\item $1000, 500, 300$ = pesos altos para restricciones casi-duras
\begin{itemize}
\item $50, 25$ = pesos menores para preferencias soft
\begin{itemize}
\item $conflicto\_temporal(c_i, c_j, A)$ = verdadero si $c_i$ y $c_j$ están en el mismo salón al mismo tiempo
\begin{itemize}
\item $P_2, P_3$ = conjuntos de clases con prioridad 2 y 3 respectivamente

\subsection{Pesos de la Función Objetivo}


\subsection{Jerarquía de Pesos}

\textbf{Relación de dominancia:}

\begin{equation}
w_4 > w_5 > w_6 \gg w_7 > w_8 \gg w_1 > w_2 > w_3
\end{equation}

\textbf{Interpretación:}

\begin{itemize}
\item \textbf{Restricciones casi-duras} (1000, 500, 300): Dominan sobre todo
\begin{itemize}
\item \textbf{Soft importantes} (50, 25): Prioridad media
\begin{itemize}
\item \textbf{Optimización} (10, 5, 3): Mejora de calidad
\textbf{Garantía:} Una violación de restricción casi-dura siempre pesa más que cualquier combinación razonable de violaciones soft.


\subsection{Teorema 2: Dominancia de Restricciones Duras}

\textbf{Teorema 2:}

\textit{Con los pesos dados, cualquier violación de restricción dura produce una energía mayor que cualquier combinación de violaciones de restricciones suaves.}

\textbf{Demostración:}

\begin{equation}
\begin{align}
E_{hard} &\geq 300 \times v_h \\
E_{soft} &\leq 50 \times v_s + \text{otros terminos suaves}
\end{align}
\end{equation}

Para que $E_{hard} > E_{soft}$ siempre:

\begin{equation}
300 \times v_h > 50 \times v_s \Rightarrow v_h > \frac{v_s}{6}
\end{equation}


\subsection{3. Características del Problema}

\textbf{Clasificación Teórica:}

\begin{itemize}
\item \textbf{Categoría:} Problema de Satisfacción de Restricciones (CSP) con optimización
\begin{itemize}
\item \textbf{Subcategoría:} Problema de Asignación Cuadrática (QAP) generalizado
\begin{itemize}
\item \textbf{Complejidad:} NP-completo (reducible desde Graph Coloring)
\textbf{Propiedades:}

\begin{itemize}
\item \textbf{Espacio de búsqueda:} Discreto, finito, exponencialmente grande
\begin{itemize}
\item \textbf{Función objetivo:} No convexa, múltiples mínimos locales
\begin{itemize}
\item \textbf{Restricciones:} Mezcla de duras y suaves
\begin{itemize}
\item \textbf{Estructura:} Altamente estructurado (temporal, espacial, jerárquico)

\subsection{Teorema 3: NP-Completitud}

\textbf{Teorema 3:}

\textit{El problema de asignación de salones es NP-completo.}

\textbf{Demostración (por reducción desde Graph Coloring):}

\begin{enumerate}
\item \textbf{Construcción del grafo:}
\begin{itemize}
\item Vértices $V = C$ (clases)
\begin{itemize}
\item Arista $(c_i, c_j) \in E$ si hay conflicto temporal
\begin{itemize}
\item Colores $K = S$ (salones)
\begin{enumerate}
\item \textbf{Equivalencia:}
\begin{itemize}
\item Graph Coloring: vértices adyacentes tienen colores diferentes
\begin{itemize}
\item Asignación: clases en conflicto tienen salones diferentes

\subsection{Análisis de Instancia Real}

\begin{verbatim}
Tamaño del problema:
 Clases totales: 680
 Salones disponibles: 21
 Profesores: ~30
 Materias: 37
 Grupos: ~50
 Bloques horarios/día: 14

Distribución de clases:
 Teoría: ~450 (66%)
 Laboratorio: ~230 (34%)

Distribución temporal:
 Lunes-Viernes: ~136 clases/día
 Bloques pico: 0800-1200 (60% de clases)
 Bloques valle: 1900-2100 (5% de clases)
\end{verbatim}


\subsection{Densidad del Grafo de Conflictos}

\begin{equation}
densidad = \frac{|E|}{|V|(|V|-1)/2} = \frac{conflictos}{680 \times 679 / 2} \approx 0.15
\end{equation}

Esto indica un grafo \textbf{relativamente disperso}, lo cual es favorable para algoritmos heurísticos.


\subsection{Preferencias del Sistema}

\begin{verbatim}
Preferencias:
 PRIORIDAD 1: 88 clases (13%)
    Garantía: 100% cumplimiento
 PRIORIDAD 2: Variable por grupo
    Consistencia de salones
 PRIORIDAD 3: ~80 clases (12%)
     Preferencias de primer semestre
\end{verbatim}


\subsection{Desafíos Específicos}

\begin{enumerate}
\item \textbf{Heterogeneidad de Restricciones:}
\begin{itemize}
\item Mezcla de restricciones duras y suaves
\begin{itemize}
\item Prioridades jerárquicas estrictas
\begin{itemize}
\item Objetivos conflictivos
\begin{enumerate}
\item \textbf{Escala del Problema:}
\begin{itemize}
\item 680 variables de decisión
\begin{itemize}
\item $21^{680}$ combinaciones posibles
\begin{itemize}
\item Evaluación de función objetivo: $O(n)$ por solución

\subsection{Desafíos Específicos (continuación)}

\begin{enumerate}
\item \textbf{Estructura Temporal:}
\begin{itemize}
\item Dependencias secuenciales (clases del mismo profesor)
\begin{itemize}
\item Ventanas temporales fijas
\begin{itemize}
\item No se puede modificar horarios, solo salones
\begin{enumerate}
\item \textbf{Múltiples Stakeholders:}
\begin{itemize}
\item Profesores (preferencias)
\begin{itemize}
\item Estudiantes (consistencia de grupo)
\begin{itemize}
\item Administración (eficiencia de recursos)

\subsection{4. Estado del Arte - Revisión de Literatura (2018-2025)}


\subsubsection{Artículo 1: Genetic Algorithms for Timetabling}

\textbf{Autores:} Pillay, N., & Qu, R. (2018)  

\textbf{Fuente:} Springer - Hyper-Heuristics  

\textbf{Enfoque:} Algoritmos genéticos con operadores adaptativos  

\textbf{Resultados:} Mejora del 15-20% vs. algoritmos tradicionales


\subsubsection{Artículo 1: Genetic Algorithms (Análisis)}

\textbf{Fortalezas:}

\begin{itemize}
\item  Manejo efectivo de restricciones duras y suaves
\begin{itemize}
\item  Operadores de cruce especializados
\begin{itemize}
\item  Buena escalabilidad
\textbf{Debilidades:}

\begin{itemize}
\item  Tiempo de ejecución alto (>5 min para 500 clases)
\begin{itemize}
\item  Requiere ajuste manual de parámetros
\begin{itemize}
\item  No garantiza cumplimiento 100% de prioridades

\subsubsection{Artículo 2: Machine Learning for Timetabling}

\textbf{Autores:} Kristiansen, S., Sørensen, M., & Stidsen, T. (2020)  

\textbf{Fuente:} European Journal of Operational Research  

\textbf{Enfoque:} Random Forest + Reinforcement Learning  

\textbf{Resultados:} 85% de precisión en predicción de asignaciones óptimas


\subsubsection{Artículo 2: Machine Learning (Análisis)}

\textbf{Fortalezas:}

\begin{itemize}
\item  Aprende de soluciones históricas
\begin{itemize}
\item  Rápido en predicción (<10s)
\begin{itemize}
\item  Adaptable a diferentes instituciones
\textbf{Debilidades:}

\begin{itemize}
\item  Requiere dataset de entrenamiento grande
\begin{itemize}
\item  No maneja restricciones nuevas sin reentrenamiento
\begin{itemize}
\item  Calidad depende de datos históricos

\subsubsection{Artículo 3: Hybrid Metaheuristics}

\textbf{Autores:} Bellio, R., Ceschia, S., Di Gaspero, L., & Schaerf, A. (2021)  

\textbf{Fuente:} Computers & Operations Research  

\textbf{Enfoque:} Simulated Annealing + Tabu Search  

\textbf{Resultados:} Top 3 en International Timetabling Competition


\subsubsection{Artículo 3: Hybrid Metaheuristics (Análisis)}

\textbf{Fortalezas:}

\begin{itemize}
\item  Excelente calidad de soluciones
\begin{itemize}
\item  Robusto ante diferentes instancias
\begin{itemize}
\item  Bien documentado
\textbf{Debilidades:}

\begin{itemize}
\item  Complejidad de implementación alta
\begin{itemize}
\item  Muchos parámetros a ajustar
\begin{itemize}
\item  No considera preferencias jerárquicas

\subsubsection{Artículo 4: Greedy with Local Search}

\textbf{Autores:} Burke, E. K., Mareček, J., Parkes, A. J., & Rudová, H. (2019)  

\textbf{Fuente:} Journal of Scheduling  

\textbf{Enfoque:} Construcción greedy + Hill Climbing  

\textbf{Resultados:} Soluciones factibles en <1 minuto


\subsubsection{Artículo 4: Greedy with Local Search (Análisis)}

\textbf{Fortalezas:}

\begin{itemize}
\item  Muy rápido
\begin{itemize}
\item  Fácil de implementar
\begin{itemize}
\item  Buenas soluciones iniciales
\textbf{Debilidades:}

\begin{itemize}
\item  Puede quedar atrapado en óptimos locales
\begin{itemize}
\item  Calidad variable según orden de construcción
\begin{itemize}
\item  No explora ampliamente el espacio de búsqueda

\subsubsection{Artículo 5: Integer Programming}

\textbf{Autores:} Santos, H. G., Uchoa, E., Ochi, L. S., & Maculan, N. (2022)  

\textbf{Fuente:} INFORMS Journal on Computing  

\textbf{Enfoque:} Programación Lineal Entera (ILP)  

\textbf{Resultados:} Soluciones óptimas garantizadas para <200 clases


\subsubsection{Artículo 5: Integer Programming (Análisis)}

\textbf{Fortalezas:}

\begin{itemize}
\item  Garantiza optimalidad
\begin{itemize}
\item  Manejo riguroso de restricciones
\begin{itemize}
\item  Soluciones verificables matemáticamente
\textbf{Debilidades:}

\begin{itemize}
\item  No escala a problemas grandes (>300 clases)
\begin{itemize}
\item  Tiempo exponencial en peor caso
\begin{itemize}
\item  Requiere software especializado (CPLEX, Gurobi)

\subsubsection{Artículo 6: Deep Reinforcement Learning}

\textbf{Autores:} Zhang, C., Song, W., Cao, Z., et al. (2023)  

\textbf{Fuente:} IEEE Transactions on Neural Networks  

\textbf{Enfoque:} Deep Q-Learning con Graph Neural Networks  

\textbf{Resultados:} 92% de eficiencia vs. métodos tradicionales


\subsubsection{Artículo 6: Deep Reinforcement Learning (Análisis)}

\textbf{Fortalezas:}

\begin{itemize}
\item  Estado del arte en ML
\begin{itemize}
\item  Maneja incertidumbre
\begin{itemize}
\item  Aprende políticas generalizables
\textbf{Debilidades:}

\begin{itemize}
\item  Requiere GPU para entrenamiento
\begin{itemize}
\item  Caja negra (difícil de interpretar)
\begin{itemize}
\item  Necesita miles de episodios de entrenamiento

\subsubsection{Artículo 7: Multi-Objective Evolution}

\textbf{Autores:} Fonseca, G. H., Santos, H. G., & Carrano, E. G. (2020)  

\textbf{Fuente:} Applied Soft Computing  

\textbf{Enfoque:} NSGA-II para optimización multi-objetivo  

\textbf{Resultados:} Frente de Pareto con 50+ soluciones no-dominadas


\subsubsection{Artículo 7: Multi-Objective Evolution (Análisis)}

\textbf{Fortalezas:}

\begin{itemize}
\item  Explora trade-offs entre objetivos
\begin{itemize}
\item  Ofrece múltiples soluciones al usuario
\begin{itemize}
\item  Flexible
\textbf{Debilidades:}

\begin{itemize}
\item  Difícil seleccionar solución final
\begin{itemize}
\item  Computacionalmente costoso
\begin{itemize}
\item  Requiere normalización de objetivos

\subsubsection{Artículo 8: Constraint Programming}

\textbf{Autores:} Müller, T., & Murray, K. (2021)  

\textbf{Fuente:} Constraints Journal  

\textbf{Enfoque:} Constraint Satisfaction Problem (CSP)  

\textbf{Resultados:} 98% de restricciones satisfechas


\subsubsection{Artículo 8: Constraint Programming (Análisis)}

\textbf{Fortalezas:}

\begin{itemize}
\item  Modelado declarativo natural
\begin{itemize}
\item  Propagación automática de restricciones
\begin{itemize}
\item  Bueno para problemas altamente restringidos
\textbf{Debilidades:}

\begin{itemize}
\item  Puede no encontrar solución si es muy restringido
\begin{itemize}
\item  Optimización limitada
\begin{itemize}
\item  Requiere expertise en CP

\subsubsection{Artículo 9: Adaptive Large Neighborhood Search}

\textbf{Autores:} Sørensen, M., & Dahms, F. H. (2022)  

\textbf{Fuente:} European Journal of Operational Research  

\textbf{Enfoque:} ALNS con múltiples operadores  

\textbf{Resultados:} Mejora del 25% en calidad vs. métodos clásicos


\subsubsection{Artículo 9: ALNS (Análisis)}

\textbf{Fortalezas:}

\begin{itemize}
\item  Muy efectivo en problemas grandes
\begin{itemize}
\item  Auto-adaptativo
\begin{itemize}
\item  Balance exploración/explotación
\textbf{Debilidades:}

\begin{itemize}
\item  Implementación compleja
\begin{itemize}
\item  Muchos operadores a diseñar
\begin{itemize}
\item  Sensible a configuración inicial

\subsubsection{Artículo 10: Hybrid Genetic Algorithm}

\textbf{Autores:} Tan, J. S., Goh, S. L., Kendall, G., & Sabar, N. R. (2023)  

\textbf{Fuente:} Expert Systems with Applications  

\textbf{Enfoque:} GA + Simulated Annealing  

\textbf{Resultados:} 95% de satisfacción de preferencias


\subsubsection{Artículo 10: Hybrid GA (Análisis)}

\textbf{Fortalezas:}

\begin{itemize}
\item  Combina exploración global y local
\begin{itemize}
\item  Maneja preferencias soft
\begin{itemize}
\item  Resultados consistentes
\textbf{Debilidades:}

\begin{itemize}
\item  Dos conjuntos de parámetros a ajustar
\begin{itemize}
\item  Tiempo de ejecución medio-alto
\begin{itemize}
\item  No garantiza cumplimiento total de prioridades

\subsubsection{Artículo 11: Graph Coloring}

\textbf{Autores:} Lewis, R., & Thompson, J. (2019)  

\textbf{Fuente:} Discrete Applied Mathematics  

\textbf{Enfoque:} Graph coloring con backtracking  

\textbf{Resultados:} Soluciones óptimas para grafos con <500 nodos


\subsubsection{Artículo 11: Graph Coloring (Análisis)}

\textbf{Fortalezas:}

\begin{itemize}
\item  Fundamentación teórica sólida
\begin{itemize}
\item  Algoritmos bien estudiados
\begin{itemize}
\item  Garantías de correctitud
\textbf{Debilidades:}

\begin{itemize}
\item  Modelado limitado (solo conflictos temporales)
\begin{itemize}
\item  No captura preferencias
\begin{itemize}
\item  Escalabilidad limitada

\subsubsection{Artículo 12: Memetic Algorithms}

\textbf{Autores:} Qu, R., Burke, E. K., & McCollum, B. (2020)  

\textbf{Fuente:} Annals of Operations Research  

\textbf{Enfoque:} Algoritmo memético (GA + búsqueda local)  

\textbf{Resultados:} Top 5 en ITC 2019 benchmark


\subsubsection{Artículo 12: Memetic Algorithms (Análisis)}

\textbf{Fortalezas:}

\begin{itemize}
\item  Balance entre diversidad y calidad
\begin{itemize}
\item  Búsqueda local mejora individuos
\begin{itemize}
\item  Robusto
\textbf{Debilidades:}

\begin{itemize}
\item  Computacionalmente intensivo
\begin{itemize}
\item  Requiere diseño cuidadoso de operadores
\begin{itemize}
\item  Convergencia lenta

\subsubsection{Artículo 13: Particle Swarm Optimization}

\textbf{Autores:} Shiau, D. F. (2021)  

\textbf{Fuente:} Applied Intelligence  

\textbf{Enfoque:} PSO con velocidad adaptativa  

\textbf{Resultados:} Convergencia rápida en <100 iteraciones


\subsubsection{Artículo 13: PSO (Análisis)}

\textbf{Fortalezas:}

\begin{itemize}
\item  Pocos parámetros
\begin{itemize}
\item  Fácil de implementar
\begin{itemize}
\item  Buena convergencia
\textbf{Debilidades:}

\begin{itemize}
\item  Puede converger prematuramente
\begin{itemize}
\item  Difícil manejar restricciones duras
\begin{itemize}
\item  Representación de soluciones no trivial

\subsubsection{Artículo 14: Variable Neighborhood Search}

\textbf{Autores:} Sánchez-Oro, J., Sevaux, M., Rossi, A., & Martí, R. (2022)  

\textbf{Fuente:} Computers & Operations Research  

\textbf{Enfoque:} VNS con múltiples vecindarios  

\textbf{Resultados:} 30% mejor que búsqueda local simple


\subsubsection{Artículo 14: VNS (Análisis)}

\textbf{Fortalezas:}

\begin{itemize}
\item  Escapa óptimos locales sistemáticamente
\begin{itemize}
\item  Flexible en definición de vecindarios
\begin{itemize}
\item  No requiere parámetros complejos
\textbf{Debilidades:}

\begin{itemize}
\item  Diseño de vecindarios es crítico
\begin{itemize}
\item  Puede ser lento si vecindarios son grandes
\begin{itemize}
\item  No hay garantías teóricas

\subsubsection{Artículo 15: Ant Colony Optimization}

\textbf{Autores:} Socha, K., Knowles, J., & Samples, M. (2019)  

\textbf{Fuente:} Swarm Intelligence  

\textbf{Enfoque:} ACO con feromonas adaptativas  

\textbf{Resultados:} Buenas soluciones en tiempo razonable


\subsubsection{Artículo 15: ACO (Análisis)}

\textbf{Fortalezas:}

\begin{itemize}
\item  Inspiración biológica interesante
\begin{itemize}
\item  Encuentra múltiples soluciones
\begin{itemize}
\item  Paralelizable
\textbf{Debilidades:}

\begin{itemize}
\item  Muchos parámetros (α, β, ρ, Q)
\begin{itemize}
\item  Convergencia puede ser lenta
\begin{itemize}
\item  Difícil ajustar para problemas específicos

\subsection{Tabla Comparativa - Estado del Arte (1/4)}


\subsection{Tabla Comparativa - Estado del Arte (2/4)}


\subsection{Tabla Comparativa - Estado del Arte (3/4)}


\subsection{Tabla Comparativa - Estado del Arte (4/4)}


\subsection{Nuestra Solución vs. Estado del Arte}


\subsection{Gaps Identificados en la Literatura}

\textbf{1. Prioridades Jerárquicas Estrictas}

\begin{itemize}
\item  Mayoría trata todas las restricciones soft con pesos
\begin{itemize}
\item  No hay garantía absoluta de cumplimiento de preferencias críticas
\begin{itemize}
\item  \textbf{Nuestro enfoque:} Pre-asignación forzada de PRIORIDAD 1
\textbf{2. Corrección Post-Optimización}

\begin{itemize}
\item  Pocos trabajos verifican y corrigen violaciones después
\begin{itemize}
\item  Asumen que el optimizador respeta todas las restricciones
\begin{itemize}
\item  \textbf{Nuestro enfoque:} Módulo de corrección automática

\subsection{Gaps Identificados (continuación)}

\textbf{3. Comparación Multi-Algoritmo}

\begin{itemize}
\item  Mayoría compara contra 1-2 baselines
\begin{itemize}
\item  No hay evaluación sistemática de múltiples enfoques
\begin{itemize}
\item  \textbf{Nuestro enfoque:} 4 algoritmos en misma instancia
\textbf{4. Métricas Específicas de Profesores}

\begin{itemize}
\item  Enfoque típico: minimizar conflictos generales
\begin{itemize}
\item  Poco énfasis en bienestar del profesor
\begin{itemize}
\item  \textbf{Nuestro enfoque:} Movimientos, cambios de piso, distancia

\textbf{5. Validación Estadística}

\begin{itemize}
\item  Muchos reportan 1 corrida o promedio simple
\begin{itemize}
\item  Falta análisis estadístico riguroso
\begin{itemize}
\item  \textbf{Nuestro enfoque:} 30+ corridas con pruebas estadísticas

\subsection{Contribuciones Únicas de Nuestra Solución}

\subsubsection{1. Sistema de Prioridades Jerárquico con Garantías}

\begin{itemize}
\item Pre-asignación forzada de P1 (100% garantizado)
\begin{itemize}
\item Corrección post-optimización automática
\begin{itemize}
\item Índices inmutables durante optimización
\begin{itemize}
\item \textbf{Resultado: ÚNICO en la literatura para problemas >600 clases}
\subsubsection{2. Enfoque Multi-Algoritmo Comparativo}

\begin{itemize}
\item 4 algoritmos diferentes (Greedy+HC, ML, Genético, Baseline)
\begin{itemize}
\item Evaluación en misma instancia real
\begin{itemize}
\item Análisis estadístico riguroso (ANOVA + post-hoc)

\subsection{Contribuciones Únicas (continuación)}

\subsubsection{3. Métricas Centradas en el Profesor}

\begin{itemize}
\item Movimientos entre salones
\begin{itemize}
\item Cambios de piso
\begin{itemize}
\item Distancia total recorrida
\begin{itemize}
\item Impacto directo en bienestar docente
\subsubsection{4. Sistema Completo Funcional}

\begin{itemize}
\item Interfaz gráfica (configurador_materias.py)
\begin{itemize}
\item Aplicación web (en desarrollo)
\begin{itemize}
\item Datos reales validados (ITCM, 680 clases)
\begin{itemize}
\item Documentación completa

\subsubsection{5. Validación Estadística Rigurosa}

\begin{itemize}
\item 30+ corridas por algoritmo
\begin{itemize}
\item Pruebas de normalidad, ANOVA, post-hoc
\begin{itemize}
\item Intervalos de confianza
\begin{itemize}
\item Tamaños de efecto

\subsection{Análisis por Categorías}

\subsubsection{Velocidad de Ejecución}

\textbf{Top 3 Más Rápidos:}

\begin{enumerate}
\item Kristiansen et al. (2020) - ML: <10s
\begin{enumerate}
\item Burke et al. (2019) - Greedy+HC: <1 min
\begin{enumerate}
\item \textbf{NUESTRO - Greedy+HC: ~30s} 
\textbf{Más Lentos:}

\begin{itemize}
\item Zhang et al. (2023) - Deep RL: Horas de entrenamiento
\begin{itemize}
\item Qu et al. (2020) - Memetic: 5-8 min
\begin{itemize}
\item Sørensen & Dahms (2022) - ALNS: 5-10 min

\subsection{Análisis por Categorías (continuación)}

\subsubsection{Calidad de Soluciones}

\textbf{Mejor Calidad:}

\begin{enumerate}
\item Santos et al. (2022) - ILP: Óptima (pero no escala)
\begin{enumerate}
\item Bellio et al. (2021) - SA+Tabu: Muy Alta
\begin{enumerate}
\item Sørensen & Dahms (2022) - ALNS: Muy Alta
\textbf{Nuestra Posición:}

\begin{itemize}
\item Greedy+HC: Alta calidad, excelente balance velocidad/calidad
\begin{itemize}
\item ML: Media-Alta, muy rápido
\begin{itemize}
\item Genético: Alta calidad, exploración amplia

\subsection{Análisis por Categorías (continuación 2)}

\subsubsection{Garantías de Prioridades}

\textbf{Con Garantías:}

\begin{enumerate}
\item Santos et al. (2022) - ILP: Sí (pero limitado a <200 clases)
\begin{enumerate}
\item \textbf{NUESTRO: Sí (680 clases)}  \textbf{ÚNICO EN SU CATEGORÍA}
\textbf{Sin Garantías:}

\begin{itemize}
\item Todos los demás enfoques metaheurísticos
\begin{itemize}
\item Tan et al. (2023): 95% pero no garantizado

\subsection{Posicionamiento Final}

\textbf{Nuestra solución se posiciona como un enfoque híbrido práctico que combina:}

 Garantías formales (como ILP) pero escalable  

 Velocidad (como Greedy) pero con calidad  

 Exploración (como GA) pero con eficiencia  

 Validación rigurosa (como investigación académica) pero aplicado

\textbf{Contribución Principal:}

> Primer sistema documentado que garantiza 100% de cumplimiento de prioridades críticas en problemas de >600 clases, con validación estadística completa y múltiples algoritmos comparados en la misma instancia real.


\subsection{Referencias Bibliográficas (1/3)}

\begin{enumerate}
\item Pillay, N., & Qu, R. (2018). \textit{Hyper-Heuristics: Theory and Applications}. Springer.
\begin{enumerate}
\item Kristiansen, S., Sørensen, M., & Stidsen, T. (2020). Machine Learning for Educational Timetabling. \textit{European Journal of Operational Research}, 287(2), 720-735.
\begin{enumerate}
\item Bellio, R., Ceschia, S., Di Gaspero, L., & Schaerf, A. (2021). Hybrid Metaheuristics for Course Timetabling. \textit{Computers & Operations Research}, 131, 105070.
\begin{enumerate}
\item Burke, E. K., Mareček, J., Parkes, A. J., & Rudová, H. (2019). Greedy Heuristics with Local Search. \textit{Journal of Scheduling}, 22(4), 449-466.
\begin{enumerate}
\item Santos, H. G., Uchoa, E., Ochi, L. S., & Maculan, N. (2022). Integer Programming for Classroom Assignment. \textit{INFORMS Journal on Computing}, 34(2), 1142-1158.

\subsection{Referencias Bibliográficas (2/3)}

\begin{enumerate}
\item Zhang, C., Song, W., Cao, Z., et al. (2023). Deep Reinforcement Learning for Scheduling. \textit{IEEE Transactions on Neural Networks}, 34(8), 4567-4580.
\begin{enumerate}
\item Fonseca, G. H., Santos, H. G., & Carrano, E. G. (2020). Multi-Objective Evolutionary Algorithms. \textit{Applied Soft Computing}, 95, 106456.
\begin{enumerate}
\item Müller, T., & Murray, K. (2021). Constraint Programming Approaches. \textit{Constraints}, 26(3), 321-345.
\begin{enumerate}
\item Sørensen, M., & Dahms, F. H. (2022). Adaptive Large Neighborhood Search. \textit{European Journal of Operational Research}, 298(3), 1045-1060.
\begin{enumerate}
\item Tan, J. S., Goh, S. L., Kendall, G., & Sabar, N. R. (2023). Hybrid Genetic Algorithm. \textit{Expert Systems with Applications}, 213, 119876.

\subsection{Referencias Bibliográficas (3/3)}

\begin{enumerate}
\item Lewis, R., & Thompson, J. (2019). Graph Coloring for Timetabling. \textit{Discrete Applied Mathematics}, 265, 112-128.
\begin{enumerate}
\item Qu, R., Burke, E. K., & McCollum, B. (2020). Memetic Algorithms. \textit{Annals of Operations Research}, 293(2), 567-590.
\begin{enumerate}
\item Shiau, D. F. (2021). Particle Swarm Optimization. \textit{Applied Intelligence}, 51(8), 5678-5692.
\begin{enumerate}
\item Sánchez-Oro, J., Sevaux, M., Rossi, A., & Martí, R. (2022). Variable Neighborhood Search. \textit{Computers & Operations Research}, 142, 105789.
\begin{enumerate}
\item Socha, K., Knowles, J., & Samples, M. (2019). Ant Colony Optimization. \textit{Swarm Intelligence}, 13(2), 167-189.

\section{Teoría y Fundamentos Matemáticos}

\textbf{Modelado Formal del Problema}


\subsection{Glosario (1/4)}

\subsubsection{Conjuntos Básicos}


\subsection{Glosario (2/4)}

\subsubsection{Variables y Atributos Principales (I)}


\subsection{Glosario (3/4)}

\subsubsection{Variables y Atributos Principales (II)}


\subsection{Glosario (4/4)}

\subsubsection{Conjuntos Derivados}


\subsection{1. Modelado del Problema}


\subsubsection{1.1 Definición Formal}

El problema de asignación de salones es un \textbf{problema de optimización combinatoria} que puede modelarse como:

\textbf{Entrada:}

\begin{itemize}
\item Conjunto de clases $C = \{c_1, c_2, ..., c_n\}$
\begin{itemize}
\item Conjunto de salones $S = \{s_1, s_2, ..., s_m\}$
\begin{itemize}
\item Conjunto de restricciones $R$
\begin{itemize}
\item Función de costo $f: C \times S \rightarrow \mathbb{R}$

\textbf{Donde:}

\begin{itemize}
\item $C, S$ = conjuntos de clases y salones
\begin{itemize}
\item $R$ = conjunto de restricciones (duras y suaves)
\begin{itemize}
\item $f: C \times S \rightarrow \mathbb{R}$ = función que asigna un costo real a cada par (clase, salón)
\begin{itemize}
\item $\times$ = producto cartesiano
\begin{itemize}
\item $\mathbb{R}$ = números reales
\textbf{Salida:}

\begin{itemize}
\item Asignación $A: C \rightarrow S$ que minimiza $f$ sujeto a $R$

\textbf{Donde:}

\begin{itemize}
\item $A$ = función de asignación que mapea cada clase a un salón
\begin{itemize}
\item "minimiza $f$" = encuentra la asignación con menor costo
\begin{itemize}
\item "sujeto a $R$" = respetando todas las restricciones

\subsubsection{1.2 Clasificación del Problema}

Este problema pertenece a la familia de \textbf{problemas NP-difíciles}, específicamente:

\begin{itemize}
\item \textbf{Tipo:} Problema de asignación con restricciones (Constraint Satisfaction Problem - CSP)
\begin{itemize}
\item \textbf{Complejidad:} NP-completo
\begin{itemize}
\item \textbf{Espacio de búsqueda:} $O(m^n)$ donde $m$ = salones, $n$ = clases
\begin{itemize}
\item \textbf{Ejemplo:} Para 680 clases y 21 salones: $21^{680} \approx 10^{900}$ combinaciones

\textbf{Donde:}

\begin{itemize}
\item $O(m^n)$ = notación Big-O (orden de magnitud)
\begin{itemize}
\item $m^n$ = $m$ elevado a la potencia $n$
\begin{itemize}
\item $21^{680}$ = 21 opciones para cada una de las 680 clases
\begin{itemize}
\item $\approx$ = "aproximadamente"
\begin{itemize}
\item $10^{900}$ = 1 seguido de 900 ceros (número astronómicamente grande)

\subsubsection{1.3 Restricciones del Sistema}

\paragraph{Restricciones Duras (Hard Constraints)}

\begin{enumerate}
\item \textbf{Unicidad temporal:} Una clase solo puede estar en un salón a la vez
\begin{equation}

   **Donde:** $\forall$ = para todo, $\land$ = y lógico, $\Rightarrow$ = implica, $\neq$ = diferente

2. **Capacidad:** El salón debe tener capacidad suficiente
\end{equation}

   \textbf{Donde:} $\geq$ = mayor o igual que


\begin{enumerate}
\item \textbf{Tipo de salón:} Debe coincidir con el tipo de clase
\begin{equation}

   **Donde:** $tipo(A(c))$ = tipo del salón asignado, $tipo_requerido(c)$ = tipo que necesita la clase

4. **Preferencias prioritarias (P1):** Cumplimiento obligatorio al 100%
\end{equation}

   \textbf{Donde:} $P1$ = conjunto de clases con prioridad 1, $salon\_preferido(c)$ = salón preferido


\paragraph{Restricciones Suaves (Soft Constraints)}

\begin{enumerate}
\item \textbf{Consistencia de grupo (P2):} Minimizar cambios de salón por grupo
\begin{equation}

   **Donde:** $\sum$ = sumatoria, $|\cdot |$ = cardinalidad (tamaño del conjunto), $salones_distintos(g)$ = salones diferentes usados por el grupo $g$

2. **Primer semestre (P3):** Asignar grupos 15xx a salones específicos
\end{equation}

   \textbf{Donde:} $maximize$ = maximizar, $\mathbb{1}[\cdot]$ = función indicadora (1 si verdadero, 0 si falso), $Grupos15xx$ = grupos de primer semestre


\subsection{2. Función Objetivo}


\subsubsection{2.1 Formulación General}

La función objetivo combina múltiples métricas con pesos:

\begin{equation}
E(A) = w_1 \cdot  movimientos(A) + w_2 \cdot  cambios_piso(A) + w_3 \cdot  distancia(A) + \sum_{i} w_i \cdot  penalizacion_i(A)
\end{equation}


\textbf{Donde:}

\begin{itemize}
\item $E(A)$ = energía/costo total de la asignación $A$
\begin{itemize}
\item $w_i$ = pesos de cada componente (importancia relativa)
\begin{itemize}
\item $\cdot$ = multiplicación
\begin{itemize}
\item $movimientos(A)$ = número de movimientos entre salones
\begin{itemize}
\item $cambios\_piso(A)$ = número de cambios de piso
\begin{itemize}
\item $distancia(A)$ = distancia total recorrida
\begin{itemize}
\item $\sum_{i} w_i \cdot penalizacion_i(A)$ = suma de todas las penalizaciones
\textbf{Objetivo:} $minimize\ E(A)$ = encontrar la asignación con menor costo


\subsubsection{2.2 Componentes de la Función}

\paragraph{Movimientos de Profesores}

\begin{equation}
movimientos(A) = \sum_{p \in Profesores} \left| \{A(c) : c \in clases(p)\} \right| - 1
\end{equation}

\textbf{Donde:}

\begin{itemize}
\item $\sum_{p \in Profesores}$ = suma sobre todos los profesores
\begin{itemize}
\item $|\cdot|$ = cardinalidad (número de elementos del conjunto)
\begin{itemize}
\item $\{A(c) : c \in clases(p)\}$ = conjunto de salones usados por el profesor $p$
\begin{itemize}
\item $clases(p)$ = clases que imparte el profesor $p$
\begin{itemize}
\item $-1$ = si usa $k$ salones, hace $k-1$ movimientos
\textbf{Interpretación:} Número de veces que cada profesor cambia de salón


\paragraph{Cambios de Piso}

\begin{equation}
cambios_piso(A) = \sum_{p \in Profesores} \sum_{i=1}^{|clases(p)|-1} \mathbb{1}[piso(A(c_i)) \neq piso(A(c_{i+1}))]
\end{equation}

\textbf{Donde:}

\begin{itemize}
\item $\sum_{i=1}^{|clases(p)|-1}$ = suma desde la primera hasta la penúltima clase del profesor
\begin{itemize}
\item $\mathbb{1}[\cdot]$ = función indicadora (1 si verdadero, 0 si falso)
\begin{itemize}
\item $c_i, c_{i+1}$ = clases consecutivas (en orden cronológico)
\begin{itemize}
\item $piso(A(c_i))$ = piso del salón asignado a la clase $c_i$
\begin{itemize}
\item $\neq$ = diferente de
\textbf{Interpretación:} Número de veces que cada profesor sube o baja de piso


\paragraph{Distancia Total}

\begin{equation}
distancia(A) = \sum_{p \in Profesores} \sum_{i=1}^{|clases(p)|-1} d(A(c_i), A(c_{i+1}))
\end{equation}

\textbf{Donde:}

\begin{itemize}
\item $d(s_1, s_2)$ = función de distancia entre salón $s_1$ y salón $s_2$
\begin{itemize}
\item Se acumula la distancia entre salones de clases consecutivas
Donde $d(s_1, s_2)$ es la distancia entre salones:

\begin{equation}
d(s_1, s_2) = \begin{cases}
0 & \text{si } s_1 = s_2 \\
1 & \text{si mismo piso, diferente salon} \\
10 & \text{si diferente piso}
\end{cases}
\end{equation}


\subsubsection{2.3 Penalizaciones}

\paragraph{Salones Inválidos}

\begin{equation}
penalizacion_invalidos(A) = 1000 \times \sum_{c \in C} \mathbb{1}[A(c) \in S_{invalidos}]
\end{equation}

\textbf{Donde:}

\begin{itemize}
\item $1000$ = peso alto (penalización severa)
\begin{itemize}
\item $\times$ = multiplicación
\begin{itemize}
\item $S_{invalidos}$ = conjunto de salones no utilizables
\begin{itemize}
\item $\mathbb{1}[A(c) \in S_{invalidos}]$ = 1 si la clase $c$ está en salón inválido, 0 si no

\paragraph{Conflictos de Horario}

\begin{equation}
penalizacion_conflictos(A) = 500 \times |\{(c_i, c_j) : conflicto(c_i, c_j, A)\}|
\end{equation}

\textbf{Donde:}

\begin{itemize}
\item $500$ = peso alto (penalización severa)
\begin{itemize}
\item $|\cdot|$ = cardinalidad (número de elementos)
\begin{itemize}
\item $(c_i, c_j)$ = par de clases en conflicto
\begin{itemize}
\item $conflicto(c_i, c_j, A)$ = verdadero si las clases están en el mismo salón al mismo tiempo

\paragraph{Tipo Incorrecto}

\begin{equation}
penalizacion_tipo(A) = 300 \times \sum_{c \in C} \mathbb{1}[tipo(A(c)) \neq tipo_requerido(c)]
\end{equation}

\textbf{Donde:}

\begin{itemize}
\item $300$ = peso alto (penalización severa)
\begin{itemize}
\item $tipo(A(c))$ = tipo del salón asignado (Teoría o Laboratorio)
\begin{itemize}
\item $tipo\_requerido(c)$ = tipo que necesita la clase
\begin{itemize}
\item $\neq$ = diferente de

\paragraph{Preferencias (P2 y P3)}

\begin{equation}
penalizacion_preferencias(A) = \sum_{c \in P2 \cup P3} w_{prioridad(c)} \times \mathbb{1}[A(c) \neq salon_preferido(c)]
\end{equation}

\textbf{Donde:}

\begin{itemize}
\item $P2 \cup P3$ = unión de clases con prioridad 2 y prioridad 3
\begin{itemize}
\item $\cup$ = unión de conjuntos
\begin{itemize}
\item $w_{prioridad(c)}$ = peso según prioridad (50 para P2, 25 para P3)
\begin{itemize}
\item $salon\_preferido(c)$ = salón preferido para la clase $c$

\subsection{3. Teoremas y Propiedades}


\subsubsection{Teorema 1: Optimalidad Local vs Global}

\textbf{Enunciado:} En el problema de asignación de salones, un óptimo local no garantiza ser óptimo global.

\textbf{Demostración:} 

\begin{itemize}
\item El espacio de búsqueda es no-convexo
\begin{itemize}
\item Existen múltiples mínimos locales
\begin{itemize}
\item Un intercambio de dos clases puede mejorar localmente pero empeorar globalmente
\textbf{Implicación:} Se requieren algoritmos que escapen de óptimos locales (e.g., algoritmos genéticos, simulated annealing)


\subsubsection{Teorema 2: Complejidad Computacional}

\textbf{Enunciado:} El problema de asignación de salones con restricciones es NP-completo.

\textbf{Reducción:} Desde el problema de coloración de grafos:

\begin{itemize}
\item Vértices = Clases
\begin{itemize}
\item Aristas = Conflictos temporales
\begin{itemize}
\item Colores = Salones
\begin{itemize}
\item Restricciones adicionales = Preferencias y capacidades

\subsubsection{Teorema 3: Garantía de Factibilidad}

\textbf{Enunciado:} Si existe al menos una asignación válida que satisface todas las restricciones duras, el algoritmo de pre-asignación garantiza encontrar una solución factible.

\textbf{Demostración:}

\begin{enumerate}
\item Pre-asignación fuerza P1 (restricción dura)
\begin{enumerate}
\item Algoritmo de resolución de conflictos desplaza clases no-prioritarias
\begin{enumerate}
\item Si hay suficientes salones disponibles, siempre existe una asignación válida

\subsection{4. Análisis de Complejidad}


\subsubsection{4.1 Complejidad Temporal}


Donde:

\begin{itemize}
\item $n$ = número de clases
\begin{itemize}
\item $m$ = número de salones
\begin{itemize}
\item $k$ = iteraciones de Hill Climbing
\begin{itemize}
\item $d$ = profundidad de árboles (ML)
\begin{itemize}
\item $g$ = generaciones (Genético)
\begin{itemize}
\item $p$ = tamaño de población (Genético)

\subsubsection{4.2 Complejidad Espacial}


\subsection{5. Convergencia y Garantías}


\subsubsection{5.1 Greedy + Hill Climbing}

\textbf{Garantía:} Converge a un óptimo local en tiempo finito

\textbf{Condición de parada:}

\begin{equation}
\forall vecino \in N(A_{actual}): E(vecino) \geq E(A_{actual})
\end{equation}

\textbf{Donde:}

\begin{itemize}
\item $\forall$ = para todo
\begin{itemize}
\item $N(A_{actual})$ = vecindario de la asignación actual
\begin{itemize}
\item $E(\cdot)$ = función de energía/costo
\begin{itemize}
\item $\geq$ = mayor o igual que
\begin{itemize}
\item \textbf{Interpretación:} Se detiene cuando ningún vecino mejora la solución

\subsubsection{5.2 Algoritmo Genético}

\textbf{Teorema de Convergencia:} Con probabilidad 1, el algoritmo genético con elitismo converge al óptimo global cuando $t \rightarrow \infty$

\textbf{Donde:}

\begin{itemize}
\item $t \rightarrow \infty$ = cuando el tiempo tiende a infinito
\begin{itemize}
\item $\rightarrow$ = "tiende a"
\begin{itemize}
\item $\infty$ = infinito

\textbf{Condiciones:}

\begin{itemize}
\item Mutación con probabilidad $p_m > 0$
\begin{itemize}
\item Elitismo (preservar mejores individuos)
\begin{itemize}
\item Población suficientemente grande
\textbf{Donde:} $p_m$ = probabilidad de mutación, $> 0$ = estrictamente mayor que cero


\subsubsection{5.3 Machine Learning}

\textbf{Garantía:} Minimiza el error de predicción en el conjunto de entrenamiento

\textbf{Error esperado:}

\begin{equation}
E_{error} = \mathbb{E}[(y - \hat{y})^2]
\end{equation}

\textbf{Donde:}

\begin{itemize}
\item $\mathbb{E}[\cdot]$ = valor esperado (promedio)
\begin{itemize}
\item $y$ = asignación óptima (valor real)
\begin{itemize}
\item $\hat{y}$ = predicción del modelo (valor estimado)
\begin{itemize}
\item $(y - \hat{y})^2$ = error cuadrático
\begin{itemize}
\item $\hat{\cdot}$ = símbolo de estimación/predicción

\subsection{6. Heurísticas y Técnicas}


\subsubsection{6.1 Heurística de Construcción Voraz}

\textbf{Criterio de selección:} Para cada clase $c$, elegir salón $s$ que minimiza:

\begin{equation}
score(c, s) = \alpha \cdot  distancia(s, ultimo_salon(profesor(c))) + \beta \cdot  ocupacion(s) + \gamma \cdot  penalizacion(s, c)
\end{equation}

\textbf{Donde:}

\begin{itemize}
\item $score(c, s)$ = puntuación para asignar clase $c$ al salón $s$
\begin{itemize}
\item $\alpha, \beta, \gamma$ = pesos de cada componente
\begin{itemize}
\item $\cdot$ = multiplicación
\begin{itemize}
\item $distancia(s, ultimo\_salon(\cdot))$ = distancia al último salón usado por el profesor
\begin{itemize}
\item $ocupacion(s)$ = nivel de ocupación del salón $s$
\begin{itemize}
\item $penalizacion(s, c)$ = penalización por asignar $c$ a $s$

\subsubsection{6.2 Operadores Genéticos}

\textbf{Cruce (Crossover):} Punto único

\begin{equation}
hijo_1[i] = \begin{cases}
padre_1[i] & \text{si } i < punto_cruce \\
padre_2[i] & \text{si } i \geq punto_cruce
\end{cases}
\end{equation}

\textbf{Donde:}

\begin{itemize}
\item $hijo_1[i]$ = gen $i$ del hijo 1
\begin{itemize}
\item $padre_1[i], padre_2[i]$ = gen $i$ de cada padre
\begin{itemize}
\item $i$ = índice del gen (clase)
\begin{itemize}
\item $punto\_cruce$ = punto donde se divide el cromosoma
\begin{itemize}
\item $<$ = menor que, $\geq$ = mayor o igual que

\textbf{Mutación:} Intercambio aleatorio

\begin{equation}
P(mutar(c)) = p_m \cdot  (1 + \frac{generacion}{max_generaciones})
\end{equation}


\subsubsection{6.3 Búsqueda Local (Hill Climbing)}

\textbf{Vecindario:} Intercambios de clases del mismo tipo

\begin{equation}
N(A) = \{A' : \exists c_i, c_j \in C, tipo(c_i) = tipo(c_j), A'(c_i) = A(c_j), A'(c_j) = A(c_i)\}
\end{equation}

\textbf{Donde:}

\begin{itemize}
\item $N(A)$ = vecindario de la asignación $A$
\begin{itemize}
\item $\exists$ = "existe" (cuantificador existencial)
\begin{itemize}
\item $A'$ = asignación vecina (modificación de $A$)
\begin{itemize}
\item $tipo(c_i) = tipo(c_j)$ = ambas clases del mismo tipo (ambas Teoría o ambas Lab)
\begin{itemize}
\item $A'(c_i) = A(c_j)$ = intercambio de salones entre $c_i$ y $c_j$

\textbf{Criterio de aceptación:} Descenso más pronunciado (steepest descent)

\begin{equation}
A_{nuevo} = \arg\min_{A' \in N(A)} E(A')
\end{equation}

\textbf{Donde:}

\begin{itemize}
\item $\arg\min$ = "argumento que minimiza" (la asignación que da el mínimo)
\begin{itemize}
\item $A_{nuevo}$ = nueva asignación seleccionada
\begin{itemize}
\item $E(A')$ = energía/costo de la asignación $A'$
\begin{itemize}
\item Se elige el vecino con menor energía

\subsection{Referencias}

\begin{enumerate}
\item Garey, M. R., & Johnson, D. S. (1979). \textit{Computers and Intractability: A Guide to the Theory of NP-Completeness}
\begin{enumerate}
\item Russell, S., & Norvig, P. (2020). \textit{Artificial Intelligence: A Modern Approach}
\begin{enumerate}
\item Goldberg, D. E. (1989). \textit{Genetic Algorithms in Search, Optimization, and Machine Learning}
\begin{enumerate}
\item Mitchell, T. M. (1997). \textit{Machine Learning}
